%% Согласно ГОСТ Р 7.0.11-2011:
%% 5.3.3 В заключении диссертации излагают итоги выполненного исследования, рекомендации, перспективы дальнейшей разработки темы.
%% 9.2.3 В заключении автореферата диссертации излагают итоги данного исследования, рекомендации и перспективы дальнейшей разработки темы.
\begin{enumerate}
  \item На основе анализа классические и нейросетевые методы восстановления траектории,
  выявить их ограничения при потерях 3--240 была разработана нейросетевая модель восстановления 
  траектории на основе селективных пространств состояний с линейной сложностью и физическим регуляризатором
  на основе закона тяги двигателей.
  \item Для дальнейшего обучения, анализа и сравнения разработанной модели были созданы синтетические в симуляторе
  Gazebo с различными динамическими характеристиками для летающих мобильных роботов в различными массо-инерциальными 
  характеристиками. 
  \item Математическое моделирование  на  синтетическом датасете показало, что разработанная нейросетевая модель обладает значительно меньшим 
  временем восстановления траектории (на участке 120 с быстрее Transformer в 80~раз) за счет линейной сложности. 
  Так же разработанная модель с регуляризатором на основе закона тяги двигателей на участке в 120 секунд показала 
  точность по среднеквадратичной ошибки траектории в 2.1 раза по сравнению с Transformer и на 42\% лучше базовой и
  на 14\% лучше обобщенной.
  \item Математическое моделирование на реальном датасете показало что разработанная нейросетевая модель восстановления 
  траектории на основе селективных пространств состояний с физическим регуляризатором на основе закона тяги двигателей показало 
  что среднеквадратичная ошибка траектории на 13\% меньше базовой.
  \item Обобщение разработанной нейронной сети на квадрокоптерах разных массы 1, 2 и 3.5 кг без дообучение показало что 
  разброс среднеквадратичная ошибка траектории на участке 3с составил 0.024м, а на 120с 0.59м. Это показывает что регуляризатор 
  на основе закона тяги позволяет обобщить результат изначальной модели на различные классы квадрокоптеров без переобучения модели. 
\end{enumerate}
